\section*{Conclusão}

Esse projeto comparou 6 algoritmos de Machine Learning na tarefa de prever o resultado de suporte em um programa de horta comunitária.
O pré-processamento foi realizando cuidadosamente, com a remoção de atributos contendo alguns dados ausentes (mais de 80\% na coluna other\_support\_plans) e o uso de codificadores ordinais e nominais foi crucial para o sucesso dos modelos.
A avaliação dos modelos foi realizada usando validação cruzada de 5 folds (5-Fold CV) para Acurácia, e uma divisão de validação (20\% dos dados) para a geração da curva AUC. Os resultados são resumidos abaixo:

\begin{center}
\label{tab:resized}
\resizebox{\textwidth}{!}{
\begin{tabular}{l c c c c c}
\toprule
\textbf{Modelo} & \textbf{Acurácia Média (CV)} & \textbf{AUC (Validação)} & \textbf{Precision} & \textbf{Recall} & \textbf{F1 score} \\
\midrule
Random Forest & 0.7437 & 0.7842 & 0.5714 & 0.2500 & 0.3478\\
Gradient Boosting & 0.7578 & 0.7440 & 0.6551 & 0.3958 & 0.4935\\
Decision Tree & 0.7031 & 0.6585 & 0.4722 & 0.3541 & 0.4047\\
Gaussian Naive Bayes & 0.7312 & 0.7539 & 0.5128 & 0.4166 & 0.4597\\
Support Vector Machine & 0.7281 & 0.6858 & 0.5416 & 0.2708 & 0.3611\\
Voting Classifier & 0.7593 & 0.7812 & 0.6800 & 0.3541 & 0.4657\\
\bottomrule
\end{tabular}
}
\end{center}

\subsection*{Análise dos Resultados}
Os resultados confirmam o poder dos métodos de ensemble (Random Forest, Gradient Boosting e Voting Classifier), que superaram significativamente a Árvore de Decisão única.
\textbf{Random Forest (0.7437)},  \textbf{Gradient Boosting (0.7578)} e \textbf{Voting Classifier (0.7593)} foram os modelos de melhor desempenho, alcançando as maiores acurácia em validação cruzada. \textbf{Random Forest (0.7842)} e \textbf{Voting Classifier (0.7812)} obtiveram as maiores AUC no conjunto de validação.

A Figura \ref{fig:roc_comp} demonstra visualmente a superioridade do Random Forest. Sua curva (AUC 0.79) está notavelmente mais próxima do canto superior esquerdo, indicativo de ser um melhor classificador em comparação com o Gradient Boosting (0.75), Naive Bayes (0.73) e a Árvore de Decisão (0.68).

\begin{figure}[H]
    \centering
    \includegraphics[width=0.5\linewidth]{results/all_models_roc_curve.png}
    \caption{Comparação da Curva ROC para todos os modelos no conjunto de validação.}
    \label{fig:roc_comp}
\end{figure}

\subsection*{Recomendação Final}
Com base nesta análise, recomenda-se à coordenação do Jardim Comunitário Raízes Vivas a implementação do modelo \textbf{Random Forest} para a tarefa de previsão.

Ele alcançou a 3ª maior acurácia em validação cruzada (0.7437)  e a maior AUC (0.7842)  no conjunto de validação, e também pode ser treinado em paralelo., demonstrando ser o modelo mais robusto e com melhor capacidade de generalização para este problema. O arquivo de submissão final \texttt{submission.csv} foi gerado pelo Random Forest.